\begin{proof}[Proof of \cref{obj:lem:cai-low-moment-duality}]
Fix a positive even integer \(K\).

\begin{enumerate}
\item Define the approximation space and target function by
\[
I_{\mathrm{appr}}=[-1,1],\qquad
\mathcal C_{\mathrm{appr}}=C(I_{\mathrm{appr}},\mathbb R),\qquad
f_{\mathrm{abs}}(x)=|x|,
\]
and equip \(\mathcal C_{\mathrm{appr}}\) with the uniform norm. Let
\[
\mathcal E_K
=\{p|_{I_{\mathrm{appr}}}:p\in\mathbb R[x],\ \deg p\le K\}.
\]
For every real polynomial \(p\),
\[
\|f_{\mathrm{abs}}-p|_{I_{\mathrm{appr}}}\|_\infty
=\sup_{|x|\le1}\bigl||x|-p(x)\bigr|.
\]
Thus \(e_K\) is the distance from \(f_{\mathrm{abs}}\) to \(\mathcal E_K\).

We first verify attainment of this distance. Define
\[
R_{\mathrm{ball}}=2\|f_{\mathrm{abs}}\|_\infty+1,\qquad
\mathcal B_K
=\{h\in\mathcal E_K:\|h\|_\infty\le R_{\mathrm{ball}}\}.
\]
The space \(\mathcal E_K\) is finite dimensional, so \(\mathcal B_K\) is compact. The continuous function
\(h\mapsto\|f_{\mathrm{abs}}-h\|_\infty\) therefore has a minimizer
\(h_{\mathrm{min}}\in\mathcal B_K\). Since \(0\in\mathcal B_K\),
\[
\|f_{\mathrm{abs}}-h_{\mathrm{min}}\|_\infty
\le\|f_{\mathrm{abs}}\|_\infty.
\]
For \(h\in\mathcal E_K\setminus\mathcal B_K\), the reverse triangle inequality gives
\[
\|f_{\mathrm{abs}}-h\|_\infty
\ge\|h\|_\infty-\|f_{\mathrm{abs}}\|_\infty
>\|f_{\mathrm{abs}}\|_\infty+1.
\]
Consequently \(h_{\mathrm{min}}\) minimizes over all of \(\mathcal E_K\), and
\[
\|f_{\mathrm{abs}}-h_{\mathrm{min}}\|_\infty=e_K.
\]
% lean: Causalean.Mathlib.Analysis.AbsoluteValueMomentPriorDuality.exists_bestPolynomialAbs

\item Let the quotient map and the target's quotient class be
\[
\pi_{\mathrm{quot}}:\mathcal C_{\mathrm{appr}}
\longrightarrow\mathcal C_{\mathrm{appr}}/\mathcal E_K,
\qquad
\bar f_{\mathrm{abs}}=\pi_{\mathrm{quot}}(f_{\mathrm{abs}}).
\]
The quotient norm is the distance to the subspace. The minimizer above gives its upper bound, while the definition of \(e_K\) gives its lower bound:
\[
\|\bar f_{\mathrm{abs}}\|
=\inf_{h\in\mathcal E_K}\|f_{\mathrm{abs}}-h\|_\infty
=e_K.
\]
By the norming-functional form of Hahn--Banach, there is a continuous real linear functional
\[
\gamma_K:\mathcal C_{\mathrm{appr}}/\mathcal E_K\longrightarrow\mathbb R,
\qquad
\|\gamma_K\|\le1,\qquad
\gamma_K(\bar f_{\mathrm{abs}})=\|\bar f_{\mathrm{abs}}\|.
\]
This also covers a zero quotient class, for which the zero functional suffices. Define
\[
\Lambda_K=\gamma_K\circ\pi_{\mathrm{quot}}.
\]
Since the quotient map is contractive and vanishes on \(\mathcal E_K\),
\[
|\Lambda_K(h)|\le\|h\|_\infty
\quad(h\in\mathcal C_{\mathrm{appr}}),\qquad
\Lambda_K(1)=0,\qquad
\Lambda_K(f_{\mathrm{abs}})=e_K.
\]
In particular, with
\[
h_\ell(x)=x^\ell
\qquad(x\in I_{\mathrm{appr}},\ 0\le\ell\le K),
\]
we have
\[
\Lambda_K(h_\ell)=0.
\]
% lean: Causalean.Mathlib.Analysis.AbsoluteValueMomentPriorDuality.exists_dual_abs_certificate

\item We construct two positive linear functionals of mass one half whose difference is \(\Lambda_K\). On
\(\mathcal C_{\mathrm{appr}}\times\mathbb R\), define the subspace, its linear functional, and a cone by
\[
\begin{aligned}
\mathcal D_K
&=\{(c_{\mathrm{const}}1,s_{\mathrm{aux}}):
       c_{\mathrm{const}},s_{\mathrm{aux}}\in\mathbb R\},\\
\beta_K(c_{\mathrm{const}}1,s_{\mathrm{aux}})
&=\frac{c_{\mathrm{const}}}{2}+s_{\mathrm{aux}},\\
\mathcal S_K
&=\{(h_++h_-,-\Lambda_K(h_-)):
       h_+,h_-\in\mathcal C_{\mathrm{appr}},\
       h_+\ge0,\ h_-\ge0\}.
\end{aligned}
\]
The representation of an element of \(\mathcal D_K\) is unique because
\(I_{\mathrm{appr}}\) is nonempty. The set \(\mathcal S_K\) is closed under addition and multiplication by nonnegative scalars.

Suppose
\[
(c_{\mathrm{const}}1,s_{\mathrm{aux}})
=(h_++h_-,-\Lambda_K(h_-))\in\mathcal D_K\cap\mathcal S_K.
\]
Then
\[
c_{\mathrm{const}}\ge0,\qquad
0\le h_-(x)\le c_{\mathrm{const}},
\]
so the function
\[
h_{\mathrm{cent}}=2h_--c_{\mathrm{const}}1
\]
satisfies
\[
\|h_{\mathrm{cent}}\|_\infty\le c_{\mathrm{const}},\qquad
2\Lambda_K(h_-)
=\Lambda_K(h_{\mathrm{cent}})
\le|\Lambda_K(h_{\mathrm{cent}})|
\le c_{\mathrm{const}}.
\]
It follows that
\[
\beta_K(c_{\mathrm{const}}1,s_{\mathrm{aux}})
=\frac{c_{\mathrm{const}}}{2}-\Lambda_K(h_-)\ge0.
\]
Moreover, for every
\((h_{\mathrm{test}},s_{\mathrm{test}})
\in\mathcal C_{\mathrm{appr}}\times\mathbb R\),
\[
(\|h_{\mathrm{test}}\|_\infty1,-s_{\mathrm{test}})
\in\mathcal D_K,
\]
and
\[
(\|h_{\mathrm{test}}\|_\infty1,-s_{\mathrm{test}})
+(h_{\mathrm{test}},s_{\mathrm{test}})
=(\|h_{\mathrm{test}}\|_\infty1+h_{\mathrm{test}},0)
\in\mathcal S_K,
\]
because the first coordinate is nonnegative. The ordered linear extension theorem therefore supplies a linear functional
\[
\Gamma_K:\mathcal C_{\mathrm{appr}}\times\mathbb R\longrightarrow\mathbb R,
\qquad
\Gamma_K|_{\mathcal D_K}=\beta_K,\qquad
\Gamma_K\ge0\ \text{on }\mathcal S_K.
\]
Define, for \(h\in\mathcal C_{\mathrm{appr}}\),
\[
\Phi_+(h)=\Gamma_K(h,0),\qquad
\Phi_-(h)=\Phi_+(h)-\Lambda_K(h).
\]
For \(h\ge0\), both \((h,0)\) and \((h,-\Lambda_K(h))\) belong to \(\mathcal S_K\). Since \(\Gamma_K(0,s)=s\) for every real \(s\),
\[
\Phi_+(h)\ge0,\qquad
\Phi_-(h)=\Gamma_K(h,-\Lambda_K(h))\ge0.
\]
Finally,
\[
\Phi_+(1)=\frac12,\qquad
\Phi_-(1)=\frac12,\qquad
\Phi_+(h)-\Phi_-(h)=\Lambda_K(h).
\]
% lean: Causalean.Mathlib.Analysis.AbsoluteValueMomentPriorDuality.exists_positive_half_split

\item Riesz representation on the compact interval gives finite positive Borel measures \(\mu_+,\mu_-\) on \(I_{\mathrm{appr}}\), characterized by
\[
\int_{I_{\mathrm{appr}}}h\,d\mu_+=\Phi_+(h),\qquad
\int_{I_{\mathrm{appr}}}h\,d\mu_-=\Phi_-(h)
\quad(h\in\mathcal C_{\mathrm{appr}}).
\]
Applying these identities to the constant function yields
\[
\mu_+(I_{\mathrm{appr}})=\mu_-(I_{\mathrm{appr}})=\frac12.
\]
Define the inclusion and reflection maps, and the symmetrized measures on \(\mathbb R\), by
\[
\iota:I_{\mathrm{appr}}\longrightarrow\mathbb R,\quad \iota(x)=x,
\qquad
r_{\mathrm{refl}}:\mathbb R\longrightarrow\mathbb R,\quad
r_{\mathrm{refl}}(x)=-x,
\]
\[
\zeta_\pm
=\frac12\left(\iota_\#\mu_\pm
 +(r_{\mathrm{refl}}\circ\iota)_\#\mu_\pm\right).
\]
Here \(\#\) denotes pushforward. Both pushforwards have mass one half and are supported on \(I_{\mathrm{appr}}\); reflection interchanges them. Hence
\[
\zeta_\pm(\mathbb R)=\frac12,\qquad
\zeta_\pm(\mathbb R\setminus I_{\mathrm{appr}})=0,\qquad
(r_{\mathrm{refl}})_\#\zeta_\pm=\zeta_\pm.
\]
For every continuous \(h_{\mathrm{test}}:\mathbb R\to\mathbb R\), the pushforward integral formula gives
\[
\int_{\mathbb R}h_{\mathrm{test}}\,d\zeta_\pm
=\frac12\left(
\int_{I_{\mathrm{appr}}}h_{\mathrm{test}}(x)\,d\mu_\pm(x)
+\int_{I_{\mathrm{appr}}}h_{\mathrm{test}}(-x)\,d\mu_\pm(x)
\right).
\]
All these integrals exist because the measures are finite and the functions are bounded on the relevant compact interval.
% lean: Causalean.Mathlib.Analysis.AbsoluteValueMomentPriorDuality.symmPush_mass
% lean: Causalean.Mathlib.Analysis.AbsoluteValueMomentPriorDuality.symmPush_supported
% lean: Causalean.Mathlib.Analysis.AbsoluteValueMomentPriorDuality.symmPush_symmetric
% lean: Causalean.Mathlib.Analysis.AbsoluteValueMomentPriorDuality.integral_symmPush

\item For \(0\le\ell\le K\), define the reflected monomial by
\[
h_{\ell,\mathrm{refl}}(x)=(-x)^\ell=(-1)^\ell h_\ell(x)
\qquad(x\in I_{\mathrm{appr}}).
\]
Linearity and polynomial annihilation give
\[
\Lambda_K(h_\ell)=0,\qquad
\Lambda_K(h_{\ell,\mathrm{refl}})
=(-1)^\ell\Lambda_K(h_\ell)=0.
\]
Combining the integral formula with the difference of the representing functionals therefore yields
\[
\begin{aligned}
\int_{\mathbb R}x^\ell\,d\zeta_+
-\int_{\mathbb R}x^\ell\,d\zeta_-
&=\frac12\left(
\Lambda_K(h_\ell)+\Lambda_K(h_{\ell,\mathrm{refl}})
\right)\\
&=0.
\end{aligned}
\]
For the absolute-value function, reflection preserves the function itself:
\[
|-x|=|x|.
\]
Consequently,
\[
\int_{\mathbb R}|x|\,d\zeta_+
-\int_{\mathbb R}|x|\,d\zeta_-
=\frac12\left(
\Lambda_K(f_{\mathrm{abs}})+\Lambda_K(f_{\mathrm{abs}})
\right)
=e_K.
\]
% lean: Causalean.Mathlib.Analysis.AbsoluteValueMomentPriorDuality.exists_absExtremalDecomposition

\item Normalize the two measures by defining
\[
\xi_0=2\zeta_-,\qquad \xi_1=2\zeta_+.
\]
Their masses are one, and multiplication by two preserves their support and reflection invariance:
\[
\xi_0(\mathbb R)=\xi_1(\mathbb R)=1,\qquad
\xi_0(\mathbb R\setminus[-1,1])
=\xi_1(\mathbb R\setminus[-1,1])=0,
\]
\[
(r_{\mathrm{refl}})_\#\xi_0=\xi_0,\qquad
(r_{\mathrm{refl}})_\#\xi_1=\xi_1.
\]
Scaling the moment identities gives, for every integer \(0\le\ell\le K\),
\[
\int x^\ell\,d\xi_1(x)-\int x^\ell\,d\xi_0(x)
=2\left(\int x^\ell\,d\zeta_+(x)-\int x^\ell\,d\zeta_-(x)\right)
=0.
\]
Likewise, scaling the absolute-moment identity gives
\[
\int |x|\,d\xi_1(x)-\int |x|\,d\xi_0(x)=2e_K.
\]
The approximation error used throughout agrees with the stated one because
\[
\inf_{\deg p\le K}\sup_{|x|\le1}\bigl||x|-p(x)\bigr|
=
\inf_{\deg p\le K}\sup_{|x|\le1}\bigl|p(x)-|x|\bigr|
=e_K.
\]
These are the required measures and identities.
% lean: Causalean.Mathlib.Analysis.AbsoluteValueMomentPriorDuality.AbsExtremalDecomposition.toMomentMatchedPriors
% lean: CausalSmith.Stat.LdpOptvalueUniformFrontier.bestAbsApproxError_eq_library
% lean: CausalSmith.Stat.LdpOptvalueUniformFrontier.caiLowMomentDuality_proved
\end{enumerate}
\end{proof}
